\section{Formula 1 Tracks}
\textbf{Note: this section is a work in progress.}

\texttt{SAT City} has been asked to host a race in the new Formula~1 season.
The mayor has tasked you with finding an optimal circuit through the city,
subject to the following requirements:
\begin{itemize}
  \item The circuit must use existing roads only, no roads may be built and no buildings destroyed.
  \item The circuit must visit at least $m$ of the $n$ points of interest in the city.
  \item The length of the circuit must be between $\ell_{\min}$ and $\ell_{\max}$.
  \item To prevent injuries, the circuit may never cross itself.
  \item The circuit must contain between $c_{\min}$ and $c_{\max}$ corners.
  \item As the cars race 50 laps, the start and end location must coincide, i.e.\ the circuit is a single closed loop.
  \item Some roads require renovation before the race may use them, due to time at most $3$ roads may be renovated.
  \item The total cost of the race (permits and renovations) must be minimized.
\end{itemize}
The mayor has provided a overview of all the roads.
He has been kind enough to color the points of interests in blue, 
and colored the roads that would need to be renovated to red.
\begin{figure}
\centering
\caption{Layout of \texttt{SAT City}}
\newcommand{\baseunit}{50}   % metres per grid square
\newcommand{\gs}{0.8}        % drawing scale: cm per grid square
% =========================================

\begin{tikzpicture}[x=0.8 cm,y=0.8 cm,
  v/.style={circle,fill,inner sep=1.4pt},
  I/.style={circle,fill,blue,inner sep=1.4pt},
  line/.style={thick,black!60!black},
  line-broken/.style={thick,red!100!black},
  lbl/.style={fill=white,inner sep=1pt,font=\tiny}]

% \edge[style]{start}{end}: draws the segment and labels its length in metres
\newcommand{\edge}[3][line]{\edgeA[#1]#2#3}
\def\edgeA[#1](#2,#3)(#4,#5){%
  \pgfmathparse{veclen((#4)-(#2),(#5)-(#3))*\baseunit}%
  \let\elen\pgfmathresult
  \draw[#1] (#2,#3) -- (#4,#5)
    node[midway,lbl] {\pgfmathprintnumber[fixed,precision=0]{\elen}};}
\newcommand{\bedge}[2]{\edge[line-broken]{#1}{#2}}

% ---- outer frame ----
\edge{(0,0)}{(2,0)}   \edge{(2,0)}{(4,0)}   \edge{(4,0)}{(6,0)}
\edge{(0,0)}{(0,-2)}  \edge{(0,-2)}{(0,-4)} \bedge{(0,-4)}{(0,-5)}
\edge{(0,-5)}{(0,-7)} \edge{(0,-7)}{(0,-10)}
\edge{(0,-10)}{(3,-10)} \edge{(3,-10)}{(5,-10)}
\edge{(6,0)}{(6,-2)}  \edge{(6,-2)}{(6,-4)} \edge{(6,-4)}{(6,-7)} \edge{(6,-7)}{(6,-8)}

% ---- top / upper middle ----
\edge{(4,0)}{(4,-0.5)} \bedge{(4,-0.5)}{(4,-3)}
\edge{(4,-3)}{(4,-5)}  \edge{(4,-5)}{(4,-6)}
\edge{(2,0)}{(2,-1)}   \edge{(2,-1)}{(2,-3)}  \edge{(2,-3)}{(2,-4)}
\edge{(3,-0.5)}{(4,-0.5)} \edge{(4,-0.5)}{(5,-0.5)} \edge{(5,-0.5)}{(6,-2)}
\edge{(3,-0.5)}{(3,-3)} \edge{(3,-3)}{(3,-4)}
\edge{(3,-0.5)}{(2,-3)}
\edge{(3,-3)}{(4,-3)} \edge{(4,-3)}{(5,-3)} \edge{(5,-3)}{(5,-4)}
\edge{(5,-4)}{(6,-4)} \edge{(6,-4)}{(7,-4)}
\edge{(3,-4)}{(2,-4)} \edge{(2,-4)}{(0,-4)}

% ---- upper left ----
\edge{(1,-1)}{(2,-1)}
\edge{(1,-1)}{(1,-2)} \edge{(1,-2)}{(1,-3)} \edge{(1,-3)}{(2,-3)}
\edge{(0,-2)}{(1,-2)}
\edge{(1,-3)}{(0,-4)}

% ---- right side ----
\edge{(6,-2)}{(8,-2)}
\edge{(8,-2)}{(8,-6)} \edge{(8,-6)}{(8,-7)}
\edge{(6,-4)}{(7,-6)}
\edge{(7,-4)}{(7,-6)} \edge{(7,-6)}{(7,-7)}
\bedge{(7,-6)}{(8,-6)}
\edge{(6,-7)}{(7,-7)} \bedge{(7,-7)}{(8,-7)}
\edge{(8,-7)}{(9,-8)} \edge{(9,-8)}{(9,-10)} \edge{(9,-10)}{(8,-11)}
\edge{(8,-11)}{(7,-12)} \edge{(7,-12)}{(4,-12)}
\edge{(8,-7)}{(8,-9)} \edge{(8,-9)}{(7,-10)} \edge{(7,-10)}{(8,-11)}
\edge{(7,-10)}{(7,-12)}
\edge{(7,-10)}{(5,-9)} \edge{(5,-9)}{(2,-9)}
\bedge{(5,-9)}{(5,-10)}

% ---- centre ----
\edge{(2,-5)}{(4,-5)}
\edge{(2,-5)}{(2,-7)} \edge{(2,-7)}{(2,-9)}
\bedge{(2,-7)}{(4,-7)}
\edge{(4,-7)}{(5,-7)} \edge{(5,-7)}{(6,-7)}
\edge{(4,-6)}{(5,-6)} \edge{(5,-6)}{(6,-7)}
\edge{(5,-6)}{(5,-7)} \edge{(5,-7)}{(5,-8)}
\edge{(4,-7)}{(4,-8)} \edge{(4,-8)}{(5,-8)} \edge{(5,-8)}{(6,-8)}
\edge{(4,-7)}{(5,-8)}

% ---- left middle ----
\edge{(0,-5)}{(0.5,-5)} \edge{(0.5,-5)}{(1,-5)}
\edge{(1,-5)}{(1,-6)} \edge{(1,-6)}{(0.5,-6)} \edge{(0.5,-6)}{(0.5,-5)}
\edge{(1,-6)}{(0,-7)} \edge{(0,-7)}{(1,-7)} \edge{(1,-7)}{(1,-6)}
\edge{(0,-7)}{(1,-8)} \edge{(1,-8)}{(0,-10)}

% ---- bottom ----
\edge{(3,-10)}{(3,-11)} \edge{(3,-11)}{(3,-12)}
\edge{(3,-12)}{(2,-12)} \edge{(2,-12)}{(2,-11)} \edge{(2,-11)}{(3,-12)}
\edge{(3,-10)}{(4,-11)} \edge{(4,-11)}{(5,-10)}
\edge{(4,-11)}{(4,-12)}
\edge{(0,-10)}{(1,-11)} \edge{(1,-11)}{(2,-12)}
\edge{(3,-12)}{(4,-12)}

% ---- vertices ----
\foreach \p in {(0,0),(2,0),(4,0),(6,0),(4,-0.5),
  (3,-0.5),(5,-0.5),(1,-1),(2,-1),
  (0,-2),(1,-2),(6,-2),(8,-2),
  (1,-3),(2,-3),(3,-3),(4,-3),(5,-3),
  (0,-4),(2,-4),(3,-4),(5,-4),(6,-4),(7,-4),
  (0,-5),(0.5,-5),(1,-5),(2,-5),(4,-5),(0.5,-6),(1,-6),
  (4,-6),(5,-6),(7,-6),(8,-6),
  (0,-7),(1,-7),(2,-7),(4,-7),(5,-7),(6,-7),(7,-7),(8,-7),
  (1,-8),(4,-8),(5,-8),(6,-8),(9,-8),
  (2,-9),(5,-9),(8,-9),
  (0,-10),(3,-10),(5,-10),(7,-10),(9,-10),
  (1,-11),(2,-11),(3,-11),(4,-11),(8,-11),
  (2,-12),(3,-12),(4,-12),(7,-12)}
  \node[v] at \p {};

% ---- "I" labels ----
\node[I] at (6,0)   {};
\node[I] at (1,-1)  {};
\node[I] at (5,-3)  {};
\node[I] at (7,-4)  {};
\node[I] at (0,-7) {};
\node[I] at (2,-7)  {};
\node[I] at (5,-10)   {};
\node[I] at (9,-10) {};
\node[I] at (3,-12) {};
\end{tikzpicture}
\end{figure}

\subsection{Solution}

\begingroup

\newcommand\ite[3]{\mathtt{ite}(#1,#2,#3)}
We identify the following parameters of the problem:
\begin{itemize}
  \item $m = 4$ is the minimum number of points of interest the track visits.
  \item $n = 6$ is the total number of points of interest in the city.
  \item $a = 7$ is the minimum number of corners.
  \item $b = 24$ is the maximum number of corners.
  \item $l = 1500$\,m is the minimum length of the track.
  \item $k = 6000$\,m is the maximum length of the track.
  \item $\delta_e \in \mathbb{N}$ is the length of road $e$.
  \item $pc_b \in \mathbb{N}$ is the permit cost of block $b$.
\end{itemize}

We use the following sets:
\begin{itemize}
  \item $V$ is the set of intersections and $E$ is the set of roads.
  \item $B$ is the set of street blocks.
  \item $E_b$ is the set of roads adjacent to block $b$.
  \item $E_v$ is the set of roads adjacent to intersection $v$.
  \item $R \subseteq E$ is the set of roads that need to be renovated.
\end{itemize}

To construct the track, we introduce the following variables:
\begin{itemize}
  \item $\alpha_e \in \mathbb{B}$ for $e \in E$: true iff road $e$ is part of the track.
  \item $\beta_{e,e'} \in \mathbb{B}$ for $e,e' \in E$: true iff the roads $e$ and $e'$ contain a point of interest.
  \item $\gamma_{e,e'} \in \mathbb{B}$ for $e,e' \in E$: true iff there is a corner between $e$ and $e'$.
  \item $p_b \in \mathbb{B}$ for $b \in B$: true iff the permit for block $b$ must be paid.
\end{itemize}


\paragraph{Cross.}
The track does not cross it self, that is each intersection has either 0 or 2 roads on the track.
For every intersection $v$ this is expressed by
\[
(\sum_{e \in E_v} \ite{\alpha_e}{1}{0} = 0 \vee \sum_{e \in V_e} \ite{\alpha_e}{1}{0} = 2)
\]

\paragraph{Points of interest.}
The track visits at least $m$ points of interest
\[
\sum_{(e,e') \in E \times E}(\ite{\alpha_{e} \wedge \alpha_{e'} \wedge \beta_{e,e'}}{1}{0}) \ge m
\]

\paragraph{Corners.}
There are between $a$ and $b$ corners
\[
a \leq \sum_{(e,e') \in E \times E} \ite{\alpha_{e} \wedge \alpha_{e'} \wedge \gamma_{e,e'}}{1}{0}) \leq b
\]

\paragraph{Renovations.}
There track requires at most $3$ roads to be renovated
\[
\sum_{e\in R } \ite{\alpha_{e} }{1}{0}) \leq 3
\]

\paragraph{Length.}
The track length is between $k$ and $l$
\[
k \leq \sum_{e \in E} \ite{\alpha_{e} }{\delta_{e}}{0} \leq l
\]

\paragraph{Permits.}
A block needs a permit if an adjacent road is part of the track.
For every block $B$ this is expressed by
\[
  \bigwedge_{e \in E_B} (\alpha_e \Rightarrow p_B).
\]

\paragraph{Objective.}
We minimize the cost of the permits and the renovations (10 per metre):
\[
 \text{minimize} \quad
 \big( \sum_{b \in B} \ite{p_b}{pc_b}{0} \big) + 
 \big( \sum_{e \in R} \ite{\alpha_e}{10 \cdot \delta_e}{0}\big)
\]



\subsection{Result}

\endgroup
